\documentclass[a4paper,11pt]{article}
\usepackage{iftex}
\ifPDFTeX
  \usepackage[utf8]{inputenc}
  \usepackage[T1]{fontenc}
  \usepackage{lmodern}
\else
  \usepackage{fontspec}
  \setmainfont{lmroman10-regular.otf}[
    BoldFont=lmroman10-bold.otf,
    ItalicFont=lmroman10-italic.otf,
    BoldItalicFont=lmroman10-bolditalic.otf]
  \setmonofont{lmmono10-regular.otf}
\fi
\usepackage[croatian]{babel}
\usepackage[margin=24mm]{geometry}
\usepackage{enumitem}
\usepackage[hidelinks]{hyperref}
\setlength{\parindent}{0pt}
\setlength{\parskip}{6pt}
\setlist[itemize]{nosep,leftmargin=16pt,topsep=3pt}
\setlist[enumerate]{nosep,leftmargin=18pt,topsep=3pt}
\hypersetup{pdftitle={Data Monitoring Platform - Izvješće 1.1},pdfauthor={Leonard Lečić i Dario Golubović}}
\AtEndDocument{\clearpage\ifnum\value{page}=3\relax\else
\PackageError{KT1}{Izvjesce mora imati dvije stranice}{Prilagodite duljinu sadrzaja.}\fi}
\begin{document}
\begin{center}
{\Large\bfseries Data Monitoring Platform}\\[5pt]
{\large Izvješće 1.1 -- Prva kontrolna točka}\\[7pt]
Leonard Lečić i Dario Golubović\\
Skriptni jezici \quad | \quad 6. studenoga 2026.
\end{center}

\section*{1. Opis problema i cilj}
Zagrebački sustav javnih bicikala Bajs putem Nextbikea objavljuje
trenutačno stanje stanica. Platforma će Python skriptom redovito
spremati ta stanja u lokalnu bazu i tako izgraditi vlastitu povijest.
Iz prikupljenih snimki utvrdit ćemo na kojim stanicama najčešće
ponestaje bicikala i u koje doba dana.

Rezultati će biti dostupni na jednoj web stranici. Korisnik će moći
odabrati stanicu s popisa i na grafu vidjeti u kojim je satima
najčešće bila prazna. Dohvat i obrada izvodit će se i kad je preglednik
zatvoren.

\section*{2. Opseg i automatizacija}
Glavni izvor je \href{https://maps.nextbike.net/maps/nextbike-live.json?city=1172&domains=hd&list_cities=0&bikes=0}{\textbf{Bajs (Nextbike Live)}}.
Osnovna verzija bilježit će identifikator i naziv stanice, broj dostupnih
bicikala, broj slobodnih mjesta kad je dostupan, status održavanja i
vrijeme dohvata. Početni interval od pet minuta moći će se promijeniti
u konfiguracijskoj datoteci. Ista skripta omogućit će ručni dohvat.

Web stranica prikazat će koliko je stanica bilo aktivno pri zadnjem
uspješnom dohvatu, koliko ih je bilo bez bicikala i kada je dohvat
obavljen. Stanice ćemo poredati prema udjelu snimki bez bicikala u
posljednjih sedam dana. Odabirom stanice prikazat će se graf po satima.

Nadogradnje ćemo birati prema raspoloživom vremenu i prikupljenim
podacima.
\begin{itemize}
\item Satni podaci o oborinama iz izvora
\href{https://open-meteo.com/en/docs/historical-weather-api}{\textbf{Open-Meteo Weather}}
za usporedbu dostupnosti u istim satima kišnih i suhih dana
\item Praćenje bicikala koji se više dana pojavljuju na istoj stanici,
provjera prekoračenja poznatog kapaciteta, klizni prosjeci i odstupanja
\item Podesivi pragovi i obavijesti e-poštom preko SMTP-a, dnevni i tjedni
sažeci te uvoz i izvoz u CSV/JSON formatu
\item Pretraga, dodatni filtri, automatsko osvježavanje prikaza i
pokretanje ručnog dohvata iz sučelja
\end{itemize}

\section*{3. Tehnologije}
Koristit ćemo \textbf{Python 3.12+}, \textbf{httpx} za dohvat,
\textbf{APScheduler} za raspored i \textbf{SQLite} uz modul
\texttt{sqlite3} za pohranu i upite. Naredbe, konfiguracija i zapis
izvršavanja oslanjat će se na module \texttt{argparse}, \texttt{tomllib}
i \texttt{logging}.

\textbf{FastAPI} i \textbf{Uvicorn} poslužit će jednu HTML stranicu i
podatke za graf. Prikaz ćemo izraditi u JavaScriptu uz \textbf{Chart.js},
a provjere dohvata i obrade pisati u \textbf{pytestu}.

\newpage
\section*{4. Arhitektura i obrada}
Jedna funkcija obavljat će dohvat, provjeru i spremanje podataka. Pozivat
će je i ručna naredba i raspoređeni posao. Web aplikacija čitat će
spremljene snimke iz iste SQLite baze.

\begin{center}
\small Dohvat $\rightarrow$ validacija $\rightarrow$ pohrana
$\rightarrow$ analiza $\rightarrow$ prikaz
\end{center}

Pri dohvatu provjerit ćemo identifikator stanice, tipove podataka i
nenegativnost broja bicikala. Vrijeme snimke spremat će se u UTC-u.
Kombinacija identifikatora stanice i vremena snimke spriječit će
dvostruki zapis iste snimke. Kod neuspjelog dohvata zabilježit ćemo
grešku, a na stranici će ostati vidljivo vrijeme posljednjeg uspjeha.

Za svaku ćemo stanicu izračunati udio valjanih snimki u kojima je broj
bicikala bio nula. Snimke stanica na održavanju izostavit ćemo iz
izračuna. Na grafu će se isti pokazatelj prikazati po lokalnom satu,
uz broj valjanih snimki za taj sat. Sati za koje nema valjanih snimki
bit će označeni kao razdoblja bez podataka.

\section*{5. Napredak u prvoj kontrolnoj točki}
Postavljen je repozitorij \texttt{data-monitoring-platform}. Nakon
povratne informacije nastavnika dogovoren je Bajs kao glavni izvor.
Definirani su osnovni opseg, tijek obrade i podjela rada. Izrada
adaptera i početak prikupljanja snimki planirani su za drugu
kontrolnu točku.

\section*{6. Plan i organizacija rada}
\begin{itemize}
\item \textbf{KT1 -- 1. studenoga.} Definicija projekta, odabir izvora
i tehnologija te plan rada.
\item \textbf{KT2 -- 27. studenoga.} Bajs adapter, ručni i raspoređeni
dohvat, SQLite pohrana i logovi. Provjera neispravnog odgovora i
dvostrukog zapisa te početak prikupljanja povijesti.
\item \textbf{KT3 -- 18. prosinca.} Izračun udjela snimki bez bicikala,
popis stanica i graf po satima. Prema mogućnostima dodati usporedbu
s oborinama ili drugu nadogradnju iz opsega.
\item \textbf{KT4 -- 22. siječnja.} Provjera cijelog postupka, ispravci
i dokumentacija. Demonstracija s prikupljenim podacima i prikazom
neuspjelog dohvata. Prema mogućnostima dodati ponovne pokušaje dohvata,
upozorenje na zastarjele podatke te sigurnosnu kopiju i oporavak baze.
\end{itemize}

Jedan član preuzet će dohvat, pohranu i raspored, a drugi analitiku
i prikaz. Oba člana radit će na Python dijelu te zajedno provjeriti
integraciju i pripremiti demonstraciju.

\end{document}

